\section*{Bab \ref{ch:b1:alcohol-activation} --- Alkohol: Mengaktifkan Gugus OH}

\begin{solution}{exo:b1:alcohol-activation:1}
\ce{2CH3CH2OH + 2Na -> 2CH3CH2O- + 2Na+ + H2};
\ce{CH3CH2OH + NaH -> CH3CH2O- + Na+ + H2};
\ce{CH3CH2OH + OH- <=> CH3CH2O- + H2O}, $K = 10^{14.0 - 15.9} = 10^{-1.9}$:
sebagian.
\end{solution}

\begin{solution}{exo:b1:alcohol-activation:2}
1-Metoksipropana; metoksietana; etoksibenzena (fenoksida, basa yang lebih
lemah daripada \omterm{def:b1:alcohol-activation:alkoxide}{alkoksida}, tetap merupakan \omterm{def:b1:electronic-effects:nucleophile}{nukleofil} yang baik).
\end{solution}

\begin{solution}{exo:b1:alcohol-activation:3}
Tosilasi dalam piridina menghasilkan \ce{CH3CH2CH2OTs}; lalu
\[ \ce{CH3CH2CH2OTs + I- -> CH3CH2CH2I + TsO-}. \] Secara langsung, iodida harus
mengusir \ce{OH-}: tidak ada reaksi. \omterm{def:b1:alcohol-activation:sulfonate}{Tosilat} mempunyai \omterm{def:b1:electronic-effects:nucleophile}{gugus pergi} yang
sangat baik, dan kondisinya tidak asam maupun sangat basa.
\end{solution}

\begin{solution}{exo:b1:alcohol-activation:4}
Butan-2-ol: but-1-ena, \cip{Z}- dan \cip{E}-but-2-ena; yang utama
\cip{E}-but-2-ena. 2-Metilpentan-2-ol: 2-metilpent-2-ena (utama,
trisubstitusi) dan 2-metilpent-1-ena.
\end{solution}

\begin{solution}{exo:b1:alcohol-activation:5}
\ce{(CH3)2CHCH2-O-CH2CH3}: natrium 2-metilpropan-1-olat dengan bromoetana,
atau natrium etoksida dengan 1-bromo-2-metilpropana. Kedua halida itu
primer, tetapi 1-bromo-2-metilpropana bercabang di sebelah karbon yang
bereaksi, sehingga SN2 melambat dan \omterm{def:b1:predominant-reaction:strong-acid}{basa kuat} itu dapat melakukan eliminasi
(2-metilpropena). Pilihan pertama lebih baik.
\end{solution}

\begin{solution}{exo:b1:alcohol-activation:6}
\omterm{def:b1:alcohol-activation:sulfonate}{Tosilat} mempertahankan \omterm{def:b1:stereochemistry-in-depth:isomers}{konfigurasi} \cip{R}; SN2 dengan etoksida
membaliknya: \cip{S}-2-etoksioktana (OEt berperingkat pertama seperti OTs
sebelumnya). Dengan asam sulfat dan kalor: eliminasi menjadi oktena
(terutama \cip{E}-okt-2-ena), melalui \omterm{def:b1:electronic-effects:intermediates}{karbokation} sekunder, dengan
hilangnya stereokimia.
\end{solution}

\begin{solution}{exo:b1:alcohol-activation:7}
Protonasi OH; lepasnya air menjadi \omterm{def:b1:electronic-effects:intermediates}{karbokation} tersier (lambat);
penangkapan oleh \ce{Cl-}: 2-kloro-2-metilpropana. Cepat karena
\omterm{def:b1:electronic-effects:intermediates}{karbokation} tersier itu stabil dan kloridanya, yang tidak larut dalam
asam, memisah.
\end{solution}

\begin{solution}{exo:b1:alcohol-activation:8}
\ce{CH3CH2OH + H+ -> CH3CH2OH2+}; molekul etanol kedua menyerang karbonnya
dari belakang sementara air pergi (SN2): \ce{(CH3CH2)2OH+}, yang
melepaskan sebuah proton: etoksietana. Yang bersaing: sebuah basa (etanol,
\ce{HSO4-}) mengambil proton $\beta$ dari \ce{CH3CH2OH2+} sementara air
pergi (E2): etena.
\end{solution}

\begin{solution}{exo:b1:alcohol-activation:9}
2-Metilpropan-2-ol (tersier), menghasilkan 2-kloro-2-metilpropana, yang tidak larut.
\end{solution}

\begin{solution}{exo:b1:alcohol-activation:10}
Protonasi dan lepasnya air menghasilkan \omterm{def:b1:electronic-effects:intermediates}{karbokation} sekunder di sebelah
karbon kuarterner. Sebuah gugus metil dari karbon itu berpindah bersama
\omterm{def:b1:lewis-resonance-vsepr:lewis}{pasangan ikatannya} ke karbon kationik: muatan positif kini berada pada
\omterm{def:b1:electronic-effects:carbon-class}{karbon tersier}, yang lebih stabil. Lepasnya proton dari karbon tetangga
menghasilkan 2,3-dimetilbut-2-ena yang tetrasubstitusi, alkena yang
paling stabil.
\end{solution}

\begin{solution}{exo:b1:alcohol-activation:11}
Dari \cip{R}-butan-2-ol: \omterm{def:b1:alcohol-activation:sulfonate}{tosilat} (\cip{R}), lalu natrium metoksida (SN2,
inversi): \cip{S}-2-metoksibutana. Dari \cip{S}-butan-2-ol: buat
\omterm{def:b1:alcohol-activation:alkoxide}{alkoksidanya} (NaH) lalu tambahkan iodometana: ikatan C--O stereogenik
tidak tersentuh, sekali lagi \cip{S}-2-metoksibutana.
\end{solution}

\begin{solution}{exo:b1:alcohol-activation:12}
$Q = [\text{eter}][\ce{H2O}]/[\ce{EtOH}]^2$. Mengeluarkan air segera
setelah terbentuk menjaga $[\ce{H2O}]$ dan $Q$ tetap kecil: $Q < K$ dan
reaksinya terus membentuk eter sampai alkoholnya habis.
\end{solution}

\begin{solution}{pb:b1:alcohol-activation:1}
\textbf{1.} Natrium tert-butoksida dengan metil halida; natrium metoksida
dengan tert-butil halida.
\textbf{2.} Pasangan kedua: metoksida, sebuah \omterm{def:b1:predominant-reaction:strong-acid}{basa kuat}, melakukan
eliminasi pada halida tersier,
\ce{(CH3)3CBr + CH3O- -> (CH3)2C=CH2 + CH3OH + Br-}.
\textbf{3.} \ce{(CH3)3CO- + CH3I -> (CH3)3COCH3 + I-}: oksigen \omterm{def:b1:alcohol-activation:alkoxide}{alkoksida}
menyerang karbon metil dari belakang sementara iodida pergi (SN2).
\textbf{4.} Dengan natrium (atau natrium hidrida): \ce{2(CH3)3COH + 2Na ->
2(CH3)3CO- + 2Na+ + H2}. Hidroksida adalah basa yang terlalu lemah: $K = 10^{14.0 -
19.2} = 10^{-5.2}$.
\textbf{5.} Air akan memprotonasi \omterm{def:b1:alcohol-activation:alkoxide}{alkoksida} dan menyerang halidanya.
\textbf{6.} Tidak: tidak ada \omterm{def:b1:stereochemistry-in-depth:stereocentre}{pusat stereogenik}.
\textbf{7.} \cip{R}-butan-2-ol + TsCl (piridina) $\to$ \cip{R}-butan-2-il
\omterm{def:b1:alcohol-activation:sulfonate}{tosilat}: retensi.
\textbf{8.} \cip{S}-2-metoksibutana, melalui inversi SN2.
\textbf{9.} \cip{S}-butan-2-ol.
\textbf{10.} \cip{R}-2-metoksibutana: ikatan C--O pada karbon stereogenik
dipertahankan; hanya O--C(metil) yang terbentuk.
\textbf{11.} E2 (metoksida adalah \omterm{def:b1:predominant-reaction:strong-acid}{basa kuat} dan karbonnya sekunder),
menghasilkan butena; dibatasi oleh basa yang kecil, \omterm{def:b1:electronic-effects:nucleophile}{gugus pergi} yang
sangat baik, dan suhu yang sedang.
\textbf{12.} Alkohol yang terprotonasi terionisasi menjadi \omterm{def:b1:electronic-effects:intermediates}{karbokation}
sekunder yang datar, yang ditangkap metanol pada kedua sisinya: eter
rasemat, bersama butena.
\textbf{13.} Protonasi OH, lepasnya air (lambat) menjadi \omterm{def:b1:electronic-effects:intermediates}{karbokation}
tersier, lepasnya proton $\beta$.
\textbf{14.} 2-Metilbut-2-ena (utama, trisubstitusi) dan
2-metilbut-1-ena.
\textbf{15.} \omterm{def:b1:electronic-effects:intermediates}{Karbokation} tersier terhalang dan jauh lebih mudah
melepaskan proton daripada diserang oleh molekul alkohol lain.
\textbf{16.} Untuk mengeluarkan sebuah produk: dengan $Q$ yang tetap kecil
reaksinya terus berlangsung, dan alkenanya tidak dibiarkan bersentuhan
dengan asam.
\textbf{17.} Sawar pertama yang tinggi (lepasnya air, penentu laju) menuju
\omterm{def:b1:electronic-effects:intermediates}{karbokation}, lalu sawar kecil menuju alkena.
\textbf{18.} Tidak: \omterm{def:b1:electronic-effects:intermediates}{karbokation} primer terlalu tidak stabil; alkohol primer
mengalami eliminasi melalui E2 pada alkohol yang terprotonasi, pada suhu
yang lebih tinggi.
\textbf{19.} $10.0/74.0 = \qty{0.135}{mol}$.
\textbf{20.} $1.10 \times 0.135 \times 141.9 = \qty{21.1}{g}$.
\textbf{21.} $0.135 \times 88.0 = \qty{11.9}{g}$.
\textbf{22.} $0.75 \times 11.9 =$ \textbf{\qty{8.9}{g} MTBE}.
\end{solution}
